\section{Introduction}
\label{sec:intro}

Airborne software loaders exchange INFORMATION and UPLOAD traffic under declared
timers, file identities and operator actions~\cite{arinc615a3}. A single missed
or late acknowledgement, a mismatched Load File, or an unresolved integrity
check can be read as a protocol failure, a measurement artefact, or an
unconfirmed injection. The engineering problem is therefore not only whether a
fixed suite eventually reports PASS, but whether a later legal action can be
chosen from residual ambiguity under a finite resource.

Requirements-based generation organises stimuli from obligations~\cite{yang2025}.
Conformance relations such as \emph{ioco} generate tests from labelled
transition systems~\cite{tretmans1996}. FSM coverage and mutant killing select
discriminating sequences~\cite{chow1978,fujiwara1991,petrenko2016,jia2011}.
Those methods address detection or selection. They do not, as stated, close an
observation--update--select loop that treats preparatory actions, unknown-effect
\texttt{ERROR}, and exclusive finite stops as first-class objects. That is a
scope comparison, not a claim that those methods cannot be extended.

This paper specifies Closed-Loop Test--Analysis Verification (CL-TAV) for
ARINC~615A data loading under timing uncertainty. The question to be answered
experimentally is: under a declared protocol, fault domain, measurement
uncertainty and resource mode, does test--analysis feedback improve verdict
reliability and isolation at an interpretable cost? Two propositions are
evaluated, not asserted as established results:
\begin{itemize}
\item C1: uncertainty-aware interpretation, timing verdicts and compatible
candidate update stay mutually consistent.
\item C2: residual ambiguity feeds the next action with explicit resource,
Prep/Recover, \texttt{ERROR} and stop rules.
\end{itemize}

The evaluation plan uses four arms (CL-T, CL-A, CL-TA, CL-LOOP) and three
registered families (DETECT, LOCATE, ABLATION). Confirmatory numbers are not
available; Section~\ref{sec:results} reserves the result layout.
\phbox{PH-ABS-RESULT}{AUTHOR-INPUT / EXPERIMENT-RESULT}{Abstract and conclusion
improvement statements remain empty until registered runs exist.}

Section~\ref{sec:related} states the problem objects.
Section~\ref{sec:arch} separates protocol CRS, tool requirements and the two
machines.
Section~\ref{sec:method} gives the loop and conditional properties.
Section~\ref{sec:setup} freezes the comparison contract.
Sections~\ref{sec:results}--\ref{sec:concl} hold results, limits and the
unfinished experimental sentence.
The supplement records ALG-CLTAV-05/06/07, wrappers and the S0--S10 map.
The paper Algorithm numbers are display labels; the stable repository IDs are
unchanged.
